\section{Introduction}


\subsection{Background and Motivation}
* Energy transition imperative: Climate goals driving decarbonization (brief)

\hrule
* Norge har historisk hatt billig hydro-kraft som har redusert behovet for thermal flexibility. Dagens situasjon, med høyere elektrisitetspriser, et stadig med constrained grid (4k mW i kø, 2.6k fra datasentre), samt variabel energiproduksjon fra fornybare kilder, er fleksibilitetsbehoved større enn aldri før. Men fjernvarmeoperatører er hindret fra å bidra til å løse dette problemet på grunn av grid tariff strukturen + får ikke deltatt i mFRR og andre fleksibilitetsmarkeder.

\subsection{Problem Statement}

\subsection{Objectives/Research Questions}

\subsection{Scope and Limitations}

\subsection{Thesis Structure}
* 1-2 sentences describing each chapter's contribution




\section{Theory}



\section{Methodology}




\section{Results}




\section{Discussion}

\subsection{Interpretation of Results}
* In depth analysis of what the results imply

\subsection{Challenges and Barriers of Implementation \& Technical vs. Market flexibility gap}

\subsection{Policy, Market and Infrastructure Implications}

\subsection{Scalability and Generalization of Results}
* Discussion of how the findings can be generalized to other parts of the grid, and how this affects the future developments

\subsection{Limitations and Future Work}
* Acknowledgment of any limitations in the methodology or data
* How might these limitations have influenced the results
* Identify areas for improvement in modeling and/or data collections
* Suggestions for further research to build on your findings




\section{Conclusion}
* Summarize the key findings and their significance
* Reflect on how the research objectives were met
* Reiterate the contributions from this thesis to the field
